
\chapter{Matching in  Observational Studies}\label{chapter::matching-obs}
 
匹配在经验研究中有很长的历史。W. Cochran 和 D. Rubin 在统计因果推断中推广了这一方法。  \citet{cochran1973controlling} 是一篇早期综述论文。\citet{rubin2006matched} 收录了 Rubin 在这一主题上的贡献。
本章还讨论 \citet{abadie2006large, abadie2008failure, abadie2011bias} 基于匹配估计量渐近分析所作的现代贡献。

 

\section{A simple starting point: many more control units}
\label{sec::ideal-matching-mpe}

\begin{figure}[ht]
\centering 
\includegraphics[width = \textwidth]{figures/matchingonetoone.pdf}
\caption{Illustration of matching in observational studies}\label{fig::matching-causal}
\end{figure}


Figure \ref{fig::matching-causal} 展示了 treatment-control observational studies 中匹配的基本思想。
考虑一个简单情形：control units 的数量 $n_0$ 远大于 treated units 的数量 $ n_1$。对 treated group 中的 unit $i=1,\ldots, n_1$，我们在 control group 中寻找一个 unit $m(i)$，使得 $X_i = X_{m(i)}$。在理想情形下，我们有精确匹配。因此，一个 matched pair 内的两个 units 有相同的 propensity score $e(X_i) = e(X_{m(i)})$。于是，在一个 unit 接受 treatment 而另一个接受 control 这一事件条件下，unit $i$ 接受 treatment 且 unit $m(i)$ 接受 control 的概率为
$$
\pr( Z_i = 1, Z_{m(i)} = 0 \mid  Z_i +Z_{m(i)} = 1, X_i, X_{m(i)} )  = 1/2
$$
这是由对称性论证得到的。\footnote{严格论证为
\begin{eqnarray*}
&&\pr( Z_i = 1, Z_{m(i)} = 0 \mid  Z_i +Z_{m(i)} = 1, X_i, X_{m(i)} ) \\
&=& \frac{ \pr( Z_i = 1, Z_{m(i)} = 0 \mid   X_i, X_{m(i)} ) }{ \pr( Z_i = 1, Z_{m(i)} = 0 \mid   X_i, X_{m(i)} )   + \pr( Z_i = 0, Z_{m(i)} = 1 \mid   X_i, X_{m(i)} ) } \\
&=& \frac{  e(X_i)  \{ 1 - e(X_{m(i)}) \}  }{ e(X_i)  \{ 1 - e(X_{m(i)}) \}   + \{1-e(X_i)\} e(X_{m(i)}) } \\
&=& \frac{1}{2}.
\end{eqnarray*}
}
也就是说，在给定 covariates 以及每一对中恰有一个 treated unit 和一个 control unit 这一事件的条件下，treatment assignment 与 MPE 完全相同。因此，我们可以把精确匹配的 observational study 当作 MPE 来分析，使用 Chapter \ref{chapter::mpe} 中的 FRT 或 Neymanian approach。这给出了关于 treated units 上 causal effect 的推断。


我们也可以为每个 treated unit 寻找多个 control units。一般地，我们可以为 treated unit $i$ 寻找 $M_i$ 个 matched control units。当 $M_i$ 可变时，这称为 {\it variable-ratio matching} \citep{ming2000substantial, ming2001note, pimentel2015variable}。在完美匹配下，treatment assignment mechanism 与 Section \ref{sec::general-matched-experiment} 中讨论的一般 matched experiment 相同。我们可以使用该节的分析结果来分析 matched observational study。


\citet{rosenbaum2002design} 提倡上述分析策略。然而，在大多数 observational studies 中，$X_i = X_{m(i)}$ 并不对所有 units 成立。上述关于 FRT 的推理不再成立。最近，\citet{guo2022statistical} 报告了 FRT 中非精确匹配后果的负面结果。这对这种策略提出了警示。
 


\section{A more complicated but realistic scenario}

即使 control group 很大，我们通常也没有精确匹配。我们能够达到的是 $X_i \approx X_{m(i)}  $，或者在某个 distance metric 下 $X_i -X_{m(i)} $ 很小。因此，我们只有近似匹配。例如，我们定义
$$
m(i) = \arg\min_{k:Z_k = 0} d(X_i ,X_k), 
$$
其中 $d(X_i ,X_k)$ 度量 $X_i$ 和 $X_k$ 之间的距离。一些标准的距离选择包括 Euclidean distance
$$
d(X_i ,X_k) =  (X_i - X_k)\tran ( X_i - X_k),
$$
以及 Mahalanobis distance
%\footnote{We define $\|v\|_2^2 = \sum_{j=1}^p v_j^2$ for a vector $v = (v_1, \ldots, v_p)\tran$. It denotes the squared length of the vector $v$.}
$$
d(X_i ,X_k) =(X_i - X_k)\tran \Omega^{-1} (X_i - X_k)
$$
其中 $\Omega$ 是来自整个 population 或仅来自 control group 的 $X_i$ 的 sample covariance matrix。


下面我回顾匹配中的一些微妙问题。综述论文见 \citet{stuart2010matching}。
\begin{enumerate}
\item
(one-to-one or one-to-$M$ matching)
上面的讨论集中在 one-to-one matching。我们也可以把讨论扩展到 one-to-$M$ matching。


\item
我关注 matching with replacement，但一些实践者更偏好 matching without replacement。如果 control units 的池子很大，这两种方法在最终结果上不会相差太多。Matching with replacement 在计算上更方便，而 matching without replacement 涉及计算量很大的 discrete optimization。Matching with replacement 通常给出质量更高的匹配，但它会因同一个 unit 被多次使用而引入依赖。相比之下，matching without replacement 的优点是 matched units 之间的独立性以及后续数据分析的简单性。


\item
由于 matched pairs 内仍有残余的 covariate imbalance，分析数据时使用 covariate adjustment 至关重要。在这种情形下，covariate adjustment 不仅是为了提高效率，也是为了 bias correction。

\item
如果 $X$ 是 ``high dimensional''，那么对 treated group 中某些 unit $i$ 以及 control group 中所有可选 units，$d(X_i ,X_k)$ 都很可能太大。如果发生这种情况，我们可能不得不丢弃一些很难找到匹配的 units。这样做实际上改变了感兴趣的 study population。


\item
上述问题很难避免。例如，如果 $X_i \sim \text{N}(0, I_p), X_k \sim\text{N}(0, I_p) ,$ 且 $X_i \ind X_k$，则
$$
(X_i - X_k)\tran (X_i - X_k) \sim    2 \chi^2_p
$$ 
其均值为 $2p$，方差为 $8p$（见 Chapter \ref{subsection::chi2-t-distribution}）。
理论表明，当 $p$ 很大时，不完美匹配会在 causal effect estimation 中造成很大的 bias。
这说明如果 $p$ 很大，我们必须在匹配之前做一些 dimension reduction。
\citet{rosenbaum1983central} 建议使用基于 propensity score 的匹配。给定 estimated propensity score，我们寻找使 $ |  \hat{e}(X_i) -  \hat{e}(X_{m(i)})  |$ 或 $| \logit\{ \hat{e}(X_i)\} - \logit\{ \hat{e}(X_{m(i)}) \} | $ 较小的 unit pairs $\{ i, m(i)\}$\footnote{定义 $\logit(w) = \log \frac{w}{1-w}$。logit function 是从 $[0,1]$ 到 $(-\infty, \infty)$ 的映射。}，也就是说，我们得到一个一维匹配问题。
\end{enumerate}

 

\section{Matching estimator for the average causal effect}

在一系列论文中，Abadie and Imbens (AI) 严格刻画了 matching estimator 的渐近性质，并为 average causal effect 提出了相应的大样本置信区间。
他们采用 observational studies 的标准设定，即 $\{ X_i , Z_i ,Y_i(1), Y_i(0)  \}_{i=1}^n \iidsim \{ X, Z, Y(1), Y(0) \}$。

\subsection{Point estimation and bias correction}

AI 关注 $1$-$M$ matching with replacement。对 treated unit $i$，我们可以简单地将 treatment 下的 potential outcome 填补为 $\hat{Y}_i(1) = Y_i$，并将 control 下的 potential outcome 填补为
$$
\hat{Y}_i(0) = M^{-1} \sum_{k\in J_i} Y_k,
$$
其中 $J_i$ 是 control group 中与 unit $i$ 匹配的 units 集合。例如，我们可以对 control group 中所有 $k$ 计算 $d(X_i, X_k)$，然后把 $J_i$ 定义为使 $d(X_i, X_k)$ 取最小 $M$ 个值的 $k$ 的指标集合。




对 control unit $i$，我们简单地将 control 下的 potential outcome 填补为 $\hat{Y}_i(0) = Y_i$，并将 treatment 下的 potential outcome 填补为
$$
\hat{Y}_i(1) = M^{-1} \sum_{k\in J_i} Y_k,
$$
其中 $J_i$ 是 treatment group 中与 unit $i$ 匹配的 units 集合。

matching estimator 为
$$
\hat{\tau}^\textup{m} = n^{-1} \sumn \{\hat{Y}_i(1) -\hat{Y}_i(0) \}.
$$

AI 表明，$\hat{\tau}^\textup{m}$ 有不可忽略的 bias，尤其是在 $X$ 是多维且 control units 的数量与 treated units 的数量相当时。通过一些技术推导，他们提出如下 bias 估计量：
$$
\hat{B} = n^{-1} \sumn  \hat{B}_i
$$
其中
$$
\hat{B}_i = (2Z_i - 1) M^{-1} \sum_{k\in J_i} \{  \hat{\mu}_{1-Z_i}(X_i) -  \hat{\mu}_{1-Z_i}(X_k)     \}
$$
其中 $\{  \hat{\mu}_1(X_i), \hat{\mu}_0(X_i) \}$ 是例如由 OLS 拟合得到的 predicted outcomes。
对 $Z_i  =1$ 的 treated unit，estimated bias 为
$$
\hat{B}_i = M^{-1} \sum_{k\in J_i} \{  \hat{\mu}_{0}(X_i) -  \hat{\mu}_{0}(X_k)     \}
$$
它校正由 covariates 不匹配导致的 predicted control potential outcomes 差异；对 $Z_i = 0$ 的 control unit，estimated bias 为
 $$
\hat{B}_i =  - M^{-1} \sum_{k\in J_i} \{  \hat{\mu}_{1}(X_i) -  \hat{\mu}_{1}(X_k)     \}
$$
它校正由 covariates 不匹配导致的 predicted treated potential outcomes 差异。


最终的 bias-corrected matching estimator 为
$$
\hat{\tau}^\textup{mbc} = \hat{\tau}^\textup{m}  - \hat{B},
$$
 它有如下 linear expansion。
 
 \begin{proposition}
 \label{prop::matching-expansion-ate} 我们有
 \begin{eqnarray}\label{eq::linearATE}
\hat{\tau}^\textup{mbc}  
= n^{-1} \sumn \hat{\psi}_i
\end{eqnarray} 
其中
 $$
  \hat{\psi}_i = \hat{\mu}_1(X_i) -  \hat{\mu}_0(X_i) + (2Z_i - 1) (1+K_i/M) \{   Y_i - \hat{\mu}_{Z_i}(X_i)  \}
 $$
 其中 $K_i$ 是 unit $i$ 被用作匹配对象的次数。
 \end{proposition}
 
 
 
Proposition  \ref{prop::matching-expansion-ate} 中的 linear expansion 来自简单但繁琐的代数。我把它的证明留作 Problem \ref{para::linear-bcm-ate}。这个 linear expansion 启发了一个简单的 variance estimator
$$
\hat{V}^\textup{mbc}  = \frac{1}{n^2} \sumn (  \hat{\psi}_i - \hat{\tau}^\textup{mbc}  )^2,
$$
这是把 $\hat{\tau}^\textup{mbc}  $ 看作 $  \hat{\psi}_i $ 的 sample average 得到的。\citet{abadie2008failure} 首先表明，通过对原始数据重抽样得到的简单 bootstrap 不能用于估计 matching estimators 的方差，但他们提出的方差估计过程并不容易实现。\citet{otsu2017bootstrap} 建议对 linear expansion 中的 $  \hat{\psi}_i $ 做 bootstrap。然而，\citet{otsu2017bootstrap} 的 bootstrap 本质上给出的就是 variance estimator $\hat{V}^\textup{mbc}  $，而这个估计量计算简单。


 
\subsection{Connection with the doubly robust estimators}
 \label{sec::connection-dr-att}
 
Bias-corrected matching estimators 和 doubly robust estimators 密切相关。二者都等于 outcome regression estimator
再加上一些基于 residuals 的修正
$$ 
 \hat{R}_i = 
 \begin{cases}
 Y_i - \hat{\mu}_1(X_i)  & \text{ if } Z_i=1;\\
 Y_i - \hat{\mu}_0(X_i)    & \text{ if } Z_i=0.
\end{cases} 
 $$
  对 average causal effect $\tau$，回忆 outcome regression estimator
$$
\hat{\tau}^{\textup{reg}} =  n^{-1} \sumn \{ \hat{\mu}_1(X_i) -  \hat{\mu}_0(X_i)\} 
$$
以及 doubly robust estimator
$$
\hat{\tau}^\textup{dr}  =  \hat{\tau}^{\textup{reg}} 
+ n^{-1} \sumn \left\{  \frac{Z_i \hat{R}_i }{ \hat{e}(X_i)} -  \frac{ (1-Z_i) \hat{R}_i}{ 1-  \hat{e}(X_i) } \right\} .
$$
进一步地，我们可以验证 $\hat{\tau}^\textup{mbc} $ 具有与 $\hat{\tau}^\textup{dr}  $ 类似的形式。
\begin{proposition}
\label{prop::mbc-dr}
$\tau$ 的 bias-corrected matching estimator 等于
$$
\hat{\tau}^\textup{mbc}   =\hat{\tau}^{\textup{reg}} 
+ n^{-1} \sumn \left\{ \left(1+{ K_i\over M} \right)  Z_i \hat{R}_i  - \left( 1+{ K_i\over M} \right) (1-Z_i)  \hat{R}_i \right\} .
$$
\end{proposition}

我把 Proposition \ref{prop::mbc-dr} 的证明留作 Problem \ref{para::dr-mbc-form}。由 Proposition \ref{prop::mbc-dr}，我们可以把 matching 看作估计 propensity score 的非参数方法，并把所得的 bias-corrected matching estimator 看作 doubly robust estimator。例如，$1+K_i /  M$ 应当接近 $1/\hat{e}(X_i)$。当一个 treated unit 的 $e(X_i)$ 很小时，基于 estimated propensity score 的权重 $1/\hat{e}(X_i)$ 会很大；与此同时，它会匹配到许多 control units，从而导致较大的 $K_i$，也就导致较大的 $1+K_i /  M$。然而，这种联系也提出了一个关于 matching 的明显问题。当 $M$ 固定时，用于 $1/e(X_i)$ 的估计量 $1+K_i /  M$ 会非常嘈杂。允许 $M$ 随 sampling size 增长，很可能改善基于 matching 的 propensity score 非参数估计量，从而改善 matching 和 bias-corrected matching estimators 的渐近性质。\citet{lin2023estimation} 给出了形式化理论，说明一旦允许 $M$ 以适当速率增长，bias-corrected matching estimator $\hat{\tau}^\textup{mbc}   $ 可以达到与 doubly robust estimator 类似的性质。


 
 
 \section{Matching estimator for the average causal effect on the treated}

 
对 average causal effect on the treated
$$
\tau_\textsc{T} = E(Y\mid Z=1) - E\{  Y(0)  \mid Z=1\},
$$ 
我们只需要为所有 treated units 填补缺失的 control 下 potential outcomes，由此得到如下估计量
$$
\hat{\tau}_\textsc{T}^\textup{m} = n_1^{-1} \sumn Z_i \{  Y_i - \hat{Y}_i(0) \}.
$$
当 $X$ 是多维时，它同样有 bias。
\citet{otsu2017bootstrap} 建议用如下方式估计它的 bias：
$$  
\hat{B}_\textsc{T} = n_1^{-1} \sumn Z_i  \hat{B}_{\textsc{T},i} 
$$
其中
$$
\hat{B}_{\textsc{T},i} = M^{-1} \sum_{k \in J_i} \{  \hat{\mu}_0(X_i) -  \hat{\mu}_0(X_k)     \}
$$
校正 $Z_i = 1$ 的 treated unit 因 covariates 不匹配而产生的 bias。


最终的 bias-corrected estimator 为
$$
\hat{\tau}_\textsc{T}^\textup{mbc} = \hat{\tau}_\textsc{T}^\textup{m} - \hat{B}_\textsc{T},
$$ 
它有如下 linear expansion。

\begin{proposition}\label{prop::linear-expansion-att}
我们有
\begin{eqnarray}\label{eq::linearATT}
\hat{\tau}_\textsc{T}^\textup{mbc}  =   n_1^{-1} \sumn \hat{\psi}_{\textsc{T}, i},
\end{eqnarray}
其中
$$
\hat{\psi}_{\textsc{T}, i} =   Z_i \{ Y_i - \hat{\mu}_0(X_i) \}
- (1-Z_i) K_i/M \{ Y_i - \hat{\mu}_0(X_i) \} . 
$$
\end{proposition}

我把证明留作 Problem \ref{para::linear-bcm-ate}。
受 \citet{otsu2017bootstrap} 启发，我们可以把 $\hat{\tau}_\textsc{T}^\textup{mbc}  $ 看作 $\hat{\psi}_{\textsc{T}, i}$ 的 sample average 乘以 $n/n_1$，因此一个直观的 variance estimator 为
$$
\hat{V}_\textsc{T}^\textup{mbc} =  \left( \frac{n}{n_1} \right)^2 \frac{1}{n^2} \sumn (\hat{\psi}_{\textsc{T}, i} - \hat{\tau}_\textsc{T}^\textup{mbc} n_1/n)^2 = \frac{1}{n_1^2} \sumn (\hat{\psi}_{\textsc{T}, i} - \hat{\tau}_\textsc{T}^\textup{mbc} n_1/n )^2.
$$

 
 类似 Section \ref{sec::connection-dr-att} 中的讨论，我们可以把 doubly robust 和 bias-corrected matching estimators 与 outcome regression estimator 作比较。
 对 average causal effect on the treated units $\tau_\textsc{T}$，回忆 outcome regression estimator
$$
\hat{\tau}^{\textup{reg}}_\textsc{T} = n_1^{-1} \sumn Z_i \{ Y_i - \hat{\mu}_0(X_i) \} ,
$$
以及 doubly robust estimator
$$
\hat{\tau}_\textsc{T}^\textup{dr} =  \hat{\tau}^{\textup{reg}}_\textsc{T} 
- n_1^{-1} \sumn \frac{  \hat{e}(X_i) }{ 1- \hat{e}(X_i) }  (1-Z_i)  \hat{R}_i .
$$
进一步地，我们可以验证 $\hat{\tau}_\textsc{T}^\textup{mbc} $ 具有与 $\hat{\tau}_\textsc{T}^\textup{dr} $ 类似的形式。
\begin{proposition}
\label{prop::mbc-dr-att}
$\tau_\textsc{T}$ 的 bias correction matching estimator 等于
$$
\hat{\tau}_\textsc{T}^\textup{mbc}   = \hat{\tau}^{\textup{reg}}_\textsc{T} 
 - n_1^{-1} \sumn { K_i\over M}  (1-Z_i)  \hat{R}_i.
$$
\end{proposition}


我把 Proposition \ref{prop::mbc-dr-att} 的证明留作 Problem \ref{para::dr-mbc-form-att}。Proposition \ref{prop::mbc-dr-att} 表明，matching 本质上使用 $K_i/M$ 来估计给定 covariates 时 treatment 的 odds。




\section{A case study}
\label{sec::matching-application}



\subsection{Experimental data}
现在我使用 \citet{sekhon2011multivariate} 的 \ri{Matching} package 重新考察 LaLonde 数据。我们已经多次对 dataset \ri{lalonde} 使用这个 package，现在将使用它的核心函数 \ri{Match}。实验部分给出如下结果：
\begin{lstlisting}
> library("car")
> library("Matching")
> 
> ## Chapter 15.5.1
> ## experimental data
> data("lalonde")
> y = lalonde$re78
> z = lalonde$treat
> x = as.matrix(lalonde[, c("age", "educ", "black",
+                           "hisp", "married", "nodegr",
+                           "re74", "re75")])
> 
> ## analysis the randomized experiment
> neymanols = lm(y ~ z)
> fisherols = lm(y ~ z + x)
> xc = scale(x)
> linols = lm(y ~ z*xc)
> resols = c(neymanols$coef[2],
+            fisherols$coef[2],
+            linols$coef[2],
+            sqrt(hccm(neymanols, type = "hc2")[2, 2]),
+            sqrt(hccm(fisherols, type = "hc2")[2, 2]),
+            sqrt(hccm(linols, type = "hc2")[2, 2]))
> resols = matrix(resols, 3, 2)
> rownames(resols) = c("neyman", "fisher", "lin")
> colnames(resols) = c("est", "se")
> resols 
            est       se
neyman 1794.343 670.9967
fisher 1676.343 677.0493
lin    1621.584 694.7217
\end{lstlisting} 
所有 regression estimators 都显示 job training program 有正的显著结果。我们可以把这些数据当作基于 1-1 matching 的 observational study 来分析，得到如下结果：
\begin{lstlisting}
> matchest.adj = Match(Y = y, Tr = z, X = x, BiasAdjust = TRUE)
> summary(matchest.adj)

Estimate...  2119.7 
AI SE......  876.42 
T-stat.....  2.4185 
p.val......  0.015583 

Original number of observations..............  445 
Original number of treated obs...............  185 
Matched number of observations...............  185 
Matched number of observations  (unweighted).  268
\end{lstlisting} 
point estimator 和 standard error 都增大了，但从定性上看，结论保持不变。


\subsection{Observational data}
接着我重新考察这组数据的 observational counterpart：
\begin{lstlisting}
> dat <- read.table("cps1re74.csv", header = TRUE)
> dat$u74 <- as.numeric(dat$re74==0)
> dat$u75 <- as.numeric(dat$re75==0)
> y = dat$re78
> z = dat$treat
> x = as.matrix(dat[, c("age", "educ", "black",
+                           "hispan", "married", "nodegree",
+                           "re74", "re75", "u74", "u75")])
 \end{lstlisting} 
如果使用简单的 OLS estimators，结果会远离 experimental benchmark，并且对 regression specification 很敏感：
\begin{lstlisting}
> neymanols = lm(y ~ z)
> fisherols = lm(y ~ z + x)
> xc = scale(x)
> linols = lm(y ~ z*xc)
> resols = c(neymanols$coef[2],
+            fisherols$coef[2],
+            linols$coef[2],
+            sqrt(hccm(neymanols, type = "hc2")[2, 2]),
+            sqrt(hccm(fisherols, type = "hc2")[2, 2]),
+            sqrt(hccm(linols, type = "hc2")[2, 2]))
> resols = matrix(resols, 3, 2)
> rownames(resols) = c("neyman", "fisher", "lin")
> colnames(resols) = c("est", "se")
> resols 
             est        se
neyman -8506.495  583.4426
fisher  1067.546  628.4389
lin    -4265.801 3211.7718
  \end{lstlisting} 
  然而，如果使用 1-1 matching，结果几乎恢复到基于 experimental data 的结果：
  \begin{lstlisting}
> matchest = Match(Y = y, Tr = z, X = x, BiasAdjust = TRUE)
> summary(matchest)

Estimate...  1747.8 
AI SE......  916.59 
T-stat.....  1.9068 
p.val......  0.056543 

Original number of observations..............  16177 
Original number of treated obs...............  185 
Matched number of observations...............  185 
Matched number of observations  (unweighted).  248
   \end{lstlisting} 
   
忽略 matched data 中的 ties，我们也可以使用 matched-pairs analysis，它同样给出与基于 experimental data 的结果相似的结果：
  \begin{lstlisting}
> diff = y[matchest$index.treated] - 
+           y[matchest$index.control]
> round(summary(lm(diff ~ 1))$coef[1, ], 2)
  Estimate Std. Error    t value   Pr(>|t|) 
   1581.44     558.55       2.83       0.01 
> 
> diff.x = x[matchest$index.treated, ] - 
+               x[matchest$index.control, ]
> round(summary(lm(diff ~ diff.x))$coef[1, ], 2)
  Estimate Std. Error    t value   Pr(>|t|) 
   1842.06     578.37       3.18       0.00 
   \end{lstlisting} 
   
   
   
   \subsection{Covariate balance checks}
   
此外，我们可以使用简单 OLS 检查 covariate balance。匹配之前，covariates 高度不平衡，这由许多 coefficients 后面的星号体现出来。
  \begin{lstlisting}
> lm.before = lm(z ~ x)
> summary(lm.before) 

Residuals:
     Min       1Q   Median       3Q      Max 
-0.18508 -0.01057  0.00303  0.01018  1.01355 

Coefficients:
              Estimate Std. Error t value Pr(>|t|)    
(Intercept)  1.404e-03  6.326e-03   0.222   0.8243    
xage        -4.043e-04  8.512e-05  -4.750 2.05e-06 ***
xeduc        3.220e-04  4.073e-04   0.790   0.4293    
xblack       1.070e-01  2.902e-03  36.871  < 2e-16 ***
xhispan      6.377e-03  3.103e-03   2.055   0.0399 *  
xmarried    -1.525e-02  2.023e-03  -7.537 5.06e-14 ***
xnodegree    1.345e-02  2.523e-03   5.331 9.89e-08 ***
xre74        7.601e-07  1.806e-07   4.208 2.59e-05 ***
xre75       -1.231e-07  1.829e-07  -0.673   0.5011    
xu74         4.224e-02  3.271e-03  12.914  < 2e-16 ***
xu75         2.424e-02  3.399e-03   7.133 1.02e-12 *** 
      \end{lstlisting} 
      
      
      然而，匹配之后，covariates 已经很好地平衡，这由所有 coefficients 都没有星号体现出来。
      
      
   \begin{lstlisting}
> lm.after = lm(z ~ x, 
+               subset = c(matchest$index.treated, 
+                          matchest$index.control))
> summary(lm.after) 

Residuals:
     Min       1Q   Median       3Q      Max 
-0.66864 -0.49161 -0.03679  0.50378  0.65122 

Coefficients:
              Estimate Std. Error t value Pr(>|t|)  
(Intercept)  6.003e-01  2.427e-01   2.474   0.0137 *
xage         3.199e-03  3.427e-03   0.933   0.3511  
xeduc       -1.501e-02  1.634e-02  -0.918   0.3590  
xblack       6.141e-05  7.408e-02   0.001   0.9993  
xhispan      1.391e-02  1.208e-01   0.115   0.9084  
xmarried    -1.328e-02  6.729e-02  -0.197   0.8437  
xnodegree   -3.023e-02  7.144e-02  -0.423   0.6723  
xre74        6.754e-06  9.864e-06   0.685   0.4939  
xre75       -9.848e-06  1.279e-05  -0.770   0.4417  
xu74         2.179e-02  1.027e-01   0.212   0.8321  
xu75        -2.642e-02  8.327e-02  -0.317   0.7512
         \end{lstlisting} 
         
         

\section{Discussion}


当 covariates 很多时，基于原始 covariates 的 matching 可能遭受 curse of dimensionality。     \citet{rosenbaum1983central} 建议使用基于 estimated propensity score 的 matching。    \citet{abadie2016matching} 为这一策略提供了形式化理论。
         
         
         
\section{Homework Problems}

\homeworksolutionref{15}
\paragraph{Linear expansions of the bias-corrected estimators}\label{para::linear-bcm-ate}


 证明 Propositions \ref{prop::matching-expansion-ate} 和 \ref{prop::linear-expansion-att}。
 
 
 
\paragraph{Doubly robust form of the bias-corrected matching estimator for $\tau$}\label{para::dr-mbc-form}

证明 Proposition \ref{prop::mbc-dr}。


\paragraph{Doubly robust form of the bias-corrected matching estimator for $\tau_\textsc{t}$}\label{para::dr-mbc-form-att}

证明 Proposition \ref{prop::mbc-dr-att}。
 

\paragraph{Revisit Example \ref{eg::chan-bmi-ATE}}


使用 matching estimator 分析 Example \ref{eg::chan-bmi-ATE} 中的 dataset。把结果与先前结果比较。
你应当检查匹配前后的 covariate balance。你也可以为 matching estimator 选择不同的匹配数量。此外，你甚至可以把各种 estimators 应用于 matched data。
你的结果对这些选择敏感吗？

 


\paragraph{Revisit Chapter \ref{sec::matching-application}}


Chapter \ref{sec::matching-application} 使用 matching 分析了 LaLonde observational study。Matching 表现很好，因为它给出的 estimator 接近 experimental gold standard。使用 outcome regression、propensity score stratification、两个 IPW 以及 doubly robust estimators 重新分析这些数据。把结果与 matching estimator 以及来自 experimental gold standard 的 estimator 比较。

注意，你有许多选择。例如，stratification 的 strata 数量，以及基于 estimated propensity scores 修剪数据的 threshold。你可以考虑拟合不同的 propensity score 和 outcome models，例如，包含 basic covariates 的一些 quadratic terms。你甚至可以把这些 estimators 应用于 matched data。

这是一个经典 dataset，已有数百篇论文使用过它。你可以阅读一些参考文献 \citep{dehejia1999causal, hainmueller2012entropy}，也可以在数据分析中发挥创造性。



\paragraph{Data re-analyses}

\citet{ho2007matching} 是 political science 中一篇有影响力的论文，作者基于它开发了一个 \ri{R} package \ri{MatchIt} \citep{ho2011matchit}。\citet{ho2007matching} 分析了两个 datasets，二者都可以从 Harvard Dataverse 获取。

使用目前讨论过的方法重新分析这两个 datasets。只要能够说明理由，你也可以尝试其他方法。



 

 
